\documentclass{article}
\usepackage{hyperref}
\usepackage[table,xcdraw]{xcolor}
\usepackage{booktabs} %Allow tables to be created well
\usepackage{graphicx}
\usepackage{float} % To make table appear where you place its code
\usepackage[margin=1in]{geometry} %border Margins
\usepackage[style=authoryear, backend=biber]{biblatex}
\addbibresource{citation.bib}


% \bibliographystyle{apa7}

\begin{document}

\title{\textbf{TOPIC:  DIET AND NUTRITION MANAGEMENT SYSTEM} }
\author{}
\date{October, 2022}
\begin{center}


\textbf{\LARGE{MAKERERE}} \includegraphics[width=70px]{muk_logo.png} \textbf{\LARGE{UNIVERSITY}}


\begin{center}
\textbf{COLLEGE OF COMPUTING AND INFORMATION SCIENCES\\ SCHOOL OF COMPUTING AND INFORMATICS TECHNOLOGY}

\end{center}

\hrule
\begin{center}
\textbf{Department of Information Systems}\\
\textbf{School of Computing and Informatics Technology}
\\ BIST:(Systems Development)
\end{center}


\vspace{15pt} 

{\LARGE TOPIC:  DIET AND NUTRITION MANAGEMENT SYSTEM}\
\vspace{9pt} 

PRESENTED BY
\vspace{-9pt} 

\begin{table}[H]
\centering
\resizebox{\columnwidth}{!}{%
\begin{tabular}{|l|l|l|l|}
\hline
\textbf{NAME} & \textbf{REG NO} & \textbf{STUDENT EMAIL ADDRESS} & \textbf{CONTACT NO} \\ \hline
MUHUMUZA VICTOR IAN & 20/U/524/PS & {\color[HTML]{000000} {victor.muhumuza@students.mak.ac.ug}} & 0761-656330 \\ \hline
WAKOKO SIMON PETER & 20/U/3512/PS & {\color[HTML]{000000} {simon.wakoko@students.mak.ac.ug}} & 0706-565245 \\ \hline
KIKOMEKO PETER GRACE & 20/U/3578/PS & {\color[HTML]{000000} { gracekikomeko@gmail.com}} & 0756-837136 \\ \hline
NAMBOOZE RACHEAL & 20/U/7994/EVE & {\color[HTML]{000000} {racheal.nambooze@students.mak.ac.ug}} & 0755-868603 \\ \hline
ONEN SAM SENSY & 20/U/3502/PS & {\color[HTML]{000000} {sensyonen@gmail.com}} & 0782-150448 \\ \hline
\end{tabular}%
}
\end{table}


\begin{center}
\textbf{A Proposal submitted to the College of Computing and Information Sciences in Partial fulfilment of the Requirements for the Award of a Degree of Bachelors of Information Systems and Technology of Makerere University.}
\end{center}

% \vspace{200pt} 
\vspace{60pt} 
\textbf{Supervisor:} Mr. Bitwire Albert George \\
\textbf{Email:} bitwire.albert@gmail.com \hspace{1cm} \textbf{Contact:} 0773-095119 \\
\textbf{Submission date:} \today

% \textbf{Option:} Systems Development
\end{center}

\newpage
%Table of content
\tableofcontents
\newpage

\section{Introduction}
\label{Introducion}
Information systems are concerned with data capture, storage, analysis, and retrieval. These are vital to assist decision-making in a short time frame, potentially allowing decisions to be made and practices to be action-ed in real-time.  In the context of diet and nutrition management, nutritional health systems allow for accurate and efficient analysis of food ingredients as well as patient nutrient intake which assists in ensuring that selected food combinations offer the desired nutrients for a particular meal or diet order (\cite{DFM}). 


\subsection{Background}
Dietary deficiency is more accountable for the overall death toll on the global scale than other factors such as tobacco, high blood pressure, or any other health risk, as indicated in the new scientific study. “Poor diet is an equal opportunity killer,” (\cite{ihme2019new}). Low levels of healthy food consumption, such as whole grains, in contrast to too many unhealthy foods, including sweetened beverages, account for one in every five deaths globally, “We are what we eat and risks affect people across a range of demographics, including age, gender, and economic status,” (\cite{ihme2019new}).

Poor diets were behind the 10.9 million deaths, or 22\% of all deaths among adults in 2017, with cardiovascular disease (CVD) as the foremost cause, trailed by cancers and diabetes. They also resulted in 255 million disability-adjusted life years (DALYs), equalling the sum of years lost and years lived with disability. Poor diet statistically represents 16\% of all DALYs among adults globally. The findings of the study dictate that while the impact of individual dietary factors varies across countries, three dietary factors – low intake of whole grains, as well as fruits, and high consumption of sodium – accounted for more than 50\% of diet-related deaths and 66\% of DALYs. The other 50\% of death and 34\% of DALYs were tied to high consumption of red meat, processed meats, sugar-sweetened beverages, and trans fatty acids among other foods. There is an urgent and compelling need for changes in the various sectors of the food production cycle, such as growing, processing, packaging, and marketing \cite{ihme2019new}.

Uganda has shown limited progress toward achieving the diet-related non-communicable disease (NCD) targets. Furthermore, it has shown no progress towards achieving the target for obesity, with an estimated 10.4\% of adult (aged 18 years and over) women and 2.3\% of adult men living with obesity. Uganda's obesity prevalence is lower than the regional average of 20.7\% for women and 9.2\% for men. At the same time, diabetes is estimated to affect 5.6\% of adult women and 5.6\% of adult men (\cite{globalnutritionreportn.d.}). The motivation for feeding should not only be to stop hunger, but also to increase the general awareness of the nutritional value of the food ingested.

\subsection{Problem statement}
Improper diet and nutrition are caused by inconsistent intake of healthy foods. Other causes include improper meal timings, under or overeating, not having enough healthy foods, and nutritional ignorance. Experts have revealed that only 10\% of children below the age of five years are eating recommended healthy foods and this includes frequenting eating nutritious meals and eating on time. This has left the majority 90\% eating non-nutritious foods which have resulted in increasing numbers of childhood malnutrition and obesity (\cite{tumwine2022only}). “People are eating a lot of unhealthy foods that don’t contribute to manufacturing of blood in the body which has made them become sicker.”(\cite{tumwine2022only}). We therefore intend to contain this problem by developing a diet and nutrition management system which will guide the targeted populace categories majorly students to maintain a healthy selection of well-balanced meals daily.

\subsection{Objectives}
\subsubsection{General Objectives}

To develop a diet and nutrition management system that will provide students with timely suggestions on the food they should consume to remain healthy. 

\subsubsection{Specific Objectives}

To identify the requirements for a Diet and Nutrition Management System in order to understand what the would-be users of the system would want.

To design a model and implement the system that will be used to assist in maintaining healthy diet meal patterns.

To test and validate the system.  

\subsection{Justification.}
Improper diet and nutrition is a constant challenge to many students which adversely exposes them to risks like obesity, digestive problems, and chronic illnesses. From the listed effects, our system intends to; 
\begin{itemize}
  \item Increase the awareness on the dire consequences of improper diet and nutrition
  \item provide a practical means to improve feeding patterns daily. 
  \item customize a system to the needs of Students Population. 

\end{itemize}

\subsection{Scope}
Our research will be focused on locally grown foods in Uganda that consist a balanced diet and specific chronic illnesses like Diabetes, obesity, pressure and cancer. We will major on Makerere University. We intend to sample several students from colleges such as CEDAT, COCIS, and CHUSS in order to know the foods they regularly eat.

\subsection{Significance}
The system we aim to build will help the elderly remind them when to best consume their meals. Athletes, through consuming the right foods will receive the required food values thus boosting their sports performance. Pregnant mothers will be able to sustain a healthy pregnancy through appropriate meals for both the unborn baby and the expectant mother. Students will be able to register an improvement in their overall mental health, especially regarding memory retention and the intelligence quotient.  People with chronic illnesses will realize the nutritional value of eating suitable foods and as a result, reduced disability-adjusted Life years. As far as academia is concerned, it will assist to refine the focus of the research on the diet and nutrition context and predict any challenges that may arise out of it.


%%-------------------------------------------------------------------------------
%% 2.0 Literature review
%%-------------------------------------------------------------------------------
\section{Literature Review}

According to Wikipedia (\citeyear{wikipedia2019dietary}), dietary management simply means providing nutritional options for individuals and groups with diet concerns through the supervision of food services. Information systems have made it easier for accurate and efficient nutritional analysis of feeding habits and patterns.

This section focuses on the exploration of different solutions developed by researchers to alleviate the problem of poor eating habits among university students. We reviewed similar systems regarding the problems faced by students when they skip meals or even fail to eat at the right time. The solution we aim to provide for this problem is to develop a diet and nutrition management system that will provide students with timely reminders, food suggestions, and nutrition literature.

We realized that so much research has been done by researchers from the field of nutrition science where a person is trained to provide information regarding the types and quantities of food the people eat. This field draws information from other areas such as biology, chemistry, and social sciences (\cite{sriram2020hire}).

The review is comprised of works that have been done to find the solution to the problem and we were able to sample various systems. We singled out systems that were closely related to our research problem and we noted the strengths and weaknesses of the solutions. By studying these systems, we were able to find the deficiency in the solutions regarding our problem and stated ways in which the solutions can be tailored to satisfy our target group. The features of the reviewed systems are compared with those of the proposed system in Table 1.

\subsection{Existing Systems}

\subsubsection{MantraCare}

MantraCare, (\cite{mantracare}), is a Ugandan web-based application that offers a range of tools and resources for individuals looking to improve their health and manage their diet. It services including consultation services with dieticians on belly fat reduction and weight loss dieting plans, exercises, and foods and may more.

\textbf{Strength of the System:}

\begin{itemize}

\item Customized meal planning: MantraCare offers personalized meal planning based on individual needs and preferences, including dietary restrictions and goals.
\item Comprehensive nutrition tracking: With MantraCare, users can track their daily intake of nutrients, including calories, carbohydrates, protein, and fats, as well as monitor their progress towards their health goals.
\item Wide range of resources: MantraCare provides access to a variety of resources, including recipe ideas, nutrition articles, and exercise guides, to help users make informed decisions about their diet and lifestyle.
\item Supportive community: MantraCare has a community of users and professionals who can provide support and guidance on nutrition and diet management.

\end{itemize}

\textbf{Weaknesses of the System:}

\begin{itemize}
\item Cost: MantraCare requires a subscription fee, which may be a barrier for some users.
\item It does not offer food suggestions.
\item Limited customization: While MantraCare offers personalized meal planning, the options may be limited for those with specific dietary needs or preferences.
\item Lack of in-person support: As an online platform, MantraCare does not offer in-person support or consultation with nutrition professionals.

\end{itemize}
\subsubsection{LifeSum}
% LifeSum is a fitness app that lets you decide what a healthy lifestyle means to you. It has plenty of recommendations and suggestions, but they are always pointing toward the goal you want.

Lifesum (\cite{lifesum}) is a digital health and wellness platform that aims to help users improve their overall health through personalized nutrition and exercise recommendations that always pointing toward the goal you want. It offers a range of features, including a food diary, macro tracking, and fitness tracking, as well as personalized meal plans and recipe suggestions.

It has several features which include:
\begin{itemize}
\item Simple tracking of meals (including barcode scanning), exercises, habits, weight, and body measurements.
\item A wide range of diets to choose between, including Keto and High Protein.
\item Favorites - save your favorite food, exercises, meals, and recipes.
\item Meal plans - 1 to 3 weeks of preplanned, easy-to-cook meals.
\item Support for macros and net carbs.
\item Detailed nutritional information.
\item Food, meal, and day ratings.
\item Weekly life score - what's gone well and how you can improve.
\item Hundreds of healthy and tasty recipes.
\item Integrates with Apple Health, Google Fit, Samsung Health, Apple Watch, RunKeeper, Fitbit, Withings, Samsung wearables, Wear OS and Google Assistant.
\end{itemize}

\textbf{Strength of the System:}
\begin{itemize}
\item The ability to integrate with different technologies for example Google Fit, and Apple Watch among others.
\item Wide range of features for comprehensive health and wellness tracking
\item Personalized meal plans and recipe suggestions
\item Integration with popular fitness tracking devices and apps
\item Available in multiple languages
\end{itemize}

\textbf{Weaknesses of the System:}
\begin{itemize}
\item Some users have reported difficulties with the food diary and tracking accuracy
\item Premium subscription required for access to certain features
\item Limited availability of certain features (e.g. personalized meal plans) in certain countries
\end{itemize}

\subsubsection{My Plate Calorie Counter}

My Plate Calorie Counter (\cite{myplate}) is a popular nutrition and diet management solution developed by the U.S. Department of Agriculture (USDA). It offers a range of features designed to help users make healthier food choices and track their daily nutrient intake. They use the world’s largest food database to provide you with calorie counts, nutritional information, and serving sizes for lots of foods. You can also access an 8-week meal plan and recipe suggestions.

\textbf{Strengths of the system:}
\begin{itemize}

\item Provides a comprehensive database of over 8,000 foods, including common brand names and restaurant meals.
\item Offers personalized daily calorie and nutrient recommendations based on the user's age, gender, weight, height, and physical activity level.
\item Allows users to track their daily intake of macronutrients (protein, carbohydrates, and fat) and micronutrients (vitamins and minerals).
\item Integrates with popular fitness apps, such as MyFitnessPal and Fitbit, to provide a more comprehensive view of the user's overall health and wellness.
\item Offers educational resources and tips on healthy eating and weight loss.

\end{itemize}

\textbf{Weaknesses of the system:}
\begin{itemize}
\item Most of the suggested foods such as Southwestern Pancakes, chocolate Almond protein cocoa, and many others are alien to our target group.
\item If you want to get most of the app such as your highest caloric foods, daily goals, and personal progress you are required to subscribe to the app which 1month costs \$9.99 (approximately 38,000/=)
\item This option would not work for our target group because the suggested foods are hard to find in Uganda and if they are available, they would be so expensive for the students.

\item Some users have reported issues with the app's tracking feature, including difficulty inputting custom foods and incorrect calorie counts.
\item Some users have reported issues with the accuracy of the nutrient information provided in the database.
\item The app is only available for iOS and Android devices, which may exclude users who do not have access to these types of devices.

\end{itemize}

\subsubsection{The Diet Planner Application}

The Diet Planner Application (\cite{dietplannerapp}), is a dietary management system that is equipped with features such as 10,000 ready meals, 270 allergens, patients' medical report capability, meal plans, personalized patient menus, nutritional interviews, kitchen and home measures, and a search engine for dishes and products, etc. To save patients' time, specialists create balanced meals that can be enjoyed without restrictions. It will allow for the preparation of nutrition plans for patients faster than ever. It enables the instant generation of fully personalized menus by a wide range of calorie goals, preferences, and dietary restrictions.

\textbf{Strength of the System:}
\begin{itemize}
\item It can generate a health report for the user.
\item It has a search engine for dishes and products.
\end{itemize}

\textbf{Weaknesses of the System:}
\begin{itemize}
\item It does not have timely reminders.
\item It does not have a BMI (Body Mass Index) calculator.
\end{itemize}

\subsubsection{Nutrition and Diet Management Solutions by Nutritics}

Nutrition and Diet Management Solutions,(\cite{nutritics}) is a software platform that provides a range of tools and resources for professionals in the field of nutrition and dietetics. The platform includes features such as a nutritional assessment tool, a meal planning module, and a recipe database, as well as educational resources such as articles and webinars. It allows multiple profile details designated as clients to be entered. Each profile is then analyzed and then, a tailored meal plan is suggested according to the profile that has been entered.

\textbf{Strengths of the system.}
\begin{itemize}
\item The Nutrition and Diet Management Solutions system offered by Nutritics includes a variety of support options for its users, including the use of webinars as a means of providing information and assistance. In addition to webinars, the system may also offer other forms of support such as email or phone support, online resources or tutorials, and in-person support through local events or workshops. This comprehensive support network helps users to make the most of the system and achieve their nutrition and diet goals.
\item It provides the user with a variety of options for meals, allowing them to choose different alternatives to fit their dietary needs and preferences.

\item It has multiple categories of users that it focuses on, including individuals looking to improve their personal nutrition and healthcare professionals looking to manage their patients' dietary needs.
\item It allows the user to add new food items and their corresponding nutritional information to the system, which helps in keeping track of the nutritional intake of individuals and making informed dietary decisions.
\item Wide range of features: Nutritics offers a range of features including meal planning, recipe analysis, and nutritional assessment tools, making it a comprehensive resource for nutrition professionals and individuals.
\item User-friendly interface: The platform has a user-friendly interface that is easy to navigate, making it accessible to users of all levels of experience.
\item Customization options: Nutritics allows users to customize their experience by setting specific goals, dietary preferences, and other parameters.
\item Accurate and reliable: Nutritics uses reliable sources and evidence-based guidelines to provide accurate and up-to-date information on nutrition and diet management.

\end{itemize}

\textbf{Weaknesses of the System.}
\begin{itemize}
\item Cost: The platform may not be accessible to all users due to the cost of subscription.
\item The application is customized for desktops only.
\item The application is customized for organizations or hospitals to manage multiple clients/patients.
\item The system is based on expensive meals that may not be affordable.
\item Lack of personalization: While Nutritics does offer customization options, it may not be able to provide the same level of personalization as a one-on-one consultation with a nutrition professional.
\end{itemize}

\subsubsection{Noom}
% Noom(\cite{noon}), is a weight loss app that uses a psychology-based approach to change your eating habits for
% the better.
Noom(\cite{noon}) is a nutrition and diet management solution that utilizes a combination of artificial intelligence and human coaching to help users make healthier lifestyle choices. One strength of Noom is its personalized approach, as the program tailors its recommendations and goals to the individual user. Noom also offers a wide range of resources and support tools, including educational articles, tracking tools, and the option for one-on-one coaching sessions


\textbf{Strengths of the system.}
\begin{itemize}
\item Noom provides a comprehensive weight loss program that includes personalized meal plans, exercise tracking, and goal setting.
\item The program also offers one-on-one coaching from certified health coaches, providing support and accountability for users.
\item Noom's focus on creating sustainable lifestyle changes rather than quick fixes is a key strength, as it helps users maintain their weight loss over the long term.
\item The program is backed by scientific research and has been shown to be effective in multiple clinical studies.
\item Noom offers a range of resources and tools for users, including a food diary, articles and tips on healthy living, and a supportive online community.
\item It is available both on mobile and desktops.
\end{itemize}

\textbf{Weaknesses of the system.}
\begin{itemize}
\item It has so many questions that take almost an hour to be answered.
\item Some users may find the cost of Noom's program to be a barrier, as it is more expensive than some other weight loss programs.
\item The program requires a significant time commitment, as users are expected to track their meals, exercise, and progress regularly.
\item Some users may find the emphasis on tracking and goal setting to be too structured, and may prefer a more flexible approach.
\item Noom's program may not be suitable for those with certain medical conditions or food allergies, as it does not offer specialized meal plans for these groups.
\item The program may not be as effective for those who are severely overweight or obese, as it is primarily focused on helping people lose smaller amounts of weight.
\item It is not accessible in terms of affordability to people like students.
\end{itemize}

\subsection{Proposed Project}
As a group, we shall develop a web application that will require university students to register possibly with their name and email addresses that will be used to send notifications to eat food cost-effectively. After this, they will be able to receive timely reminders on the minimum amount and type of food they should consume to successfully fulfill their daily tasks. This will include a mix of foods such as Matooke, Rice, and Sweet potatoes and drinks such as water and a cup of tea. It will also include general knowledge on how they can maintain a healthy and functioning body such as drinking lots of water, exercising regularly, and avoiding harmful consumption of alcohol, and drugs among others.
This will provide a more customized and cost-effective feeding way for university students to allow them to focus on their studies as they also take care of their health.

\begin{table}[H]
\centering
\caption{Comparison Of the Systems}
\label{tab:comparison-of -systems}
\resizebox{\columnwidth}{!}{%
\begin{tabular}{@{}|l|l|l|l|l|l|l|l|@{}}
\toprule
\multicolumn{1}{|c|}{\textbf{System Features}} & \textbf{MantraCare} & \textbf{\begin{tabular}[c]{@{}l@{}}Nutrition and Diet \\ Management Solutions\\  By Nutritics\end{tabular}} & \textbf{Noom} & \textbf{\begin{tabular}[c]{@{}l@{}}Diet \\ Planner\end{tabular}} & \textbf{\begin{tabular}[c]{@{}l@{}}My Plate \\ Calories\end{tabular}} & \textbf{Life Sum} & \textbf{\begin{tabular}[c]{@{}l@{}}Proposed \\ System\end{tabular}} \\ \midrule
BMI Calculator &  & \textbf{YES} & \textbf{YES} &  & \textbf{YES} & \textbf{YES} & \textbf{YES} \\ \midrule
Timely Reminder & \textbf{YES} & \textbf{YES} &  &  & \textbf{YES} &  & \textbf{YES} \\ \midrule
Food Suggestions & \textbf{YES} & \textbf{YES} & \textbf{YES} & \textbf{YES} &  &  & \textbf{YES} \\ \midrule
Nutritional Literature & \textbf{YES} &  & \textbf{YES} & \textbf{YES} & \textbf{YES} & \textbf{YES} & \textbf{YES} \\ \midrule
Accessibility/Affordable(Free) &  &  &  &  &  &  & \textbf{YES} \\ \bottomrule
\end{tabular}%
}
\end{table}

% \bibliographystyle{apa}% Set the bibliography style
% \bibliography{citation} % Include the bibliography file
\newpage
\printbibliography

\end{document}
